\documentclass[12pt]{article}
\usepackage[T1]{fontenc}
\usepackage[utf8]{inputenc}
\usepackage{lmodern}
\usepackage{textcomp}
\usepackage{enumitem}
\usepackage{geometry}
\geometry{margin=1in}
\begin{document}
\begin{center}
Answer Key - Computer Science - Undergraduate Level Assessment 13 \
\small (Answers are ordered exactly as in the question paper.)
\end{center}

\section*{Multiple Choice}
\begin{enumerate}[label=\arabic*.]
\item \textbf{Answer:} B
\\ \textit{Options:} A) 2; B) 6; C) 3; D) 1
\\ \textit{Note:} The correct answer is 6 because the lambda function r takes an input q and multiplies it by 2, so when r is called with the argument 3, it computes 3 * 2, resulting in 6.
\item \textbf{Answer:} C
\\ \textit{Options:} A) It is used to call procedures with addresses that are farther than $2^16$ bytes away.; B) It cannot return a value.; C) It cannot pass parameters by reference.; D) It cannot call procedures implemented in a different language.
\\ \textit{Note:} Remote procedure calls (RPC) typically serialize parameters for transmission over a network, which means they are passed by value rather than by reference, preventing direct manipulation of the original parameter data on the caller's side.
\item \textbf{Answer:} A
\\ \textit{Options:} A) isupper(); B) join(seq); C) len(string); D) ljust(width[, fillchar])
\\ \textit{Note:} The isupper() function in Python 3 checks if all the characters in a string are uppercase, returning True if they are and False otherwise, making it the correct choice for determining if a string is entirely in uppercase.
\item \textbf{Answer:} D
\\ \textit{Options:} A) 1; B) 6; C) 4; D) 5
\\ \textit{Note:} The correct answer is 5 because in Python, the modulo operator (\%) has a higher precedence than addition (+), so the expression evaluates as 4 + (3 \% 2), which simplifies to 4 + 1, resulting in 5.
\item \textbf{Answer:} B
\\ \textit{Options:} A) /; B) //; C) \%; D) |
\\ \textit{Note:} The correct answer is "//" because in Python 3, floor division is performed using this operator, which divides two numbers and rounds down the result to the nearest whole number.
\end{enumerate}

\section*{True / False}
\begin{enumerate}[label=\arabic*.]
\setcounter{enumi}{5}
\item \textbf{Answer:} True
\\ \textit{Options:} A) True; B) False
\\ \textit{Note:} In Python, variable names are case-sensitive, which means that the programming language differentiates between uppercase and lowercase letters, treating 'Variable' and 'variable' as two separate identifiers.
\item \textbf{Answer:} True
\\ \textit{Options:} A) True; B) False
\\ \textit{Note:} In a complete K-ary tree of depth N, the total number of nodes is given by the formula ($K^(N$+1) - 1) / (K - 1), while the number of nonterminal (internal) nodes is $K^N$ - 1, leading to a ratio of nonterminal nodes to total nodes that approaches 1/K as N increases.
\item \textbf{Answer:} False
\\ \textit{Options:} A) True; B) False
\\ \textit{Note:} The minimum value in the list l = [1,2,3,4] is 1, not 2, because the minimum function identifies the smallest element in the list.
\item \textbf{Answer:} False
\\ \textit{Options:} A) True; B) False
\\ \textit{Note:} The string 'xxzy' cannot be generated by the provided grammar because it does not conform to the production rules defined by the grammar that dictate the structure and composition of valid strings.
\item \textbf{Answer:} True
\\ \textit{Options:} A) True; B) False
\\ \textit{Note:} The expression evaluates to 1 because the modulo operation 3 \% 3 results in 0, so the overall expression 1 + 0 equals 1.
\end{enumerate}

\section*{Long Form Response}
\begin{enumerate}[label=\arabic*.]
\setcounter{enumi}{10}
\item \textbf{Answer:} def div\_list(nums 1,nums 2): | return list(result)
\\ \textit{Note:} correct because it demonstrates the use of a lambda function with the map function to perform element-wise division between two lists, thereby achieving the desired result of dividing corresponding elements from the input lists.
\item \textbf{Answer:} def sort\_tuple(tup): | return tup
\\ \textit{Note:} incorrect because it does not implement any sorting logic for the list of tuples based on the last element of each tuple.
\end{enumerate}

\vfill
\noindent\textit{End of Answer Key}
\end{document}